\PassOptionsToPackage{unicode=true}{hyperref} % options for packages loaded elsewhere
\PassOptionsToPackage{hyphens}{url}
%
\documentclass[]{article}
\usepackage{lmodern}
\usepackage{amssymb,amsmath}
\usepackage{ifxetex,ifluatex}
\usepackage{fixltx2e} % provides \textsubscript
\ifnum 0\ifxetex 1\fi\ifluatex 1\fi=0 % if pdftex
  \usepackage[T1]{fontenc}
  \usepackage[utf8]{inputenc}
  \usepackage{textcomp} % provides euro and other symbols
\else % if luatex or xelatex
  \usepackage{unicode-math}
  \defaultfontfeatures{Ligatures=TeX,Scale=MatchLowercase}
\fi
% use upquote if available, for straight quotes in verbatim environments
\IfFileExists{upquote.sty}{\usepackage{upquote}}{}
% use microtype if available
\IfFileExists{microtype.sty}{%
\usepackage[]{microtype}
\UseMicrotypeSet[protrusion]{basicmath} % disable protrusion for tt fonts
}{}
\IfFileExists{parskip.sty}{%
\usepackage{parskip}
}{% else
\setlength{\parindent}{0pt}
\setlength{\parskip}{6pt plus 2pt minus 1pt}
}
\usepackage{hyperref}
\hypersetup{
            pdfborder={0 0 0},
            breaklinks=true}
\urlstyle{same}  % don't use monospace font for urls
\usepackage{longtable,booktabs}
% Fix footnotes in tables (requires footnote package)
\IfFileExists{footnote.sty}{\usepackage{footnote}\makesavenoteenv{longtable}}{}
\usepackage{graphicx,grffile}
\makeatletter
\def\maxwidth{\ifdim\Gin@nat@width>\linewidth\linewidth\else\Gin@nat@width\fi}
\def\maxheight{\ifdim\Gin@nat@height>\textheight\textheight\else\Gin@nat@height\fi}
\makeatother
% Scale images if necessary, so that they will not overflow the page
% margins by default, and it is still possible to overwrite the defaults
% using explicit options in \includegraphics[width, height, ...]{}
\setkeys{Gin}{width=\maxwidth,height=\maxheight,keepaspectratio}
\setlength{\emergencystretch}{3em}  % prevent overfull lines
\providecommand{\tightlist}{%
  \setlength{\itemsep}{0pt}\setlength{\parskip}{0pt}}
\setcounter{secnumdepth}{0}
% Redefines (sub)paragraphs to behave more like sections
\ifx\paragraph\undefined\else
\let\oldparagraph\paragraph
\renewcommand{\paragraph}[1]{\oldparagraph{#1}\mbox{}}
\fi
\ifx\subparagraph\undefined\else
\let\oldsubparagraph\subparagraph
\renewcommand{\subparagraph}[1]{\oldsubparagraph{#1}\mbox{}}
\fi

% set default figure placement to htbp
\makeatletter
\def\fps@figure{htbp}
\makeatother


\date{}

\begin{document}

\textbf{From Expansion to Compression: Adaptive Multi-View
Recommendation with LLM-based Semantic Denoising}

\textbf{ABSTRACT}

Graph-based Collaborative Filtering (CF) and Large Language Models
(LLMs) represent two distinct paradigms in recommendation, yet they
suffer from complementary limitations: CF models are semantically blind
to content-driven user intents, while LLMs lack population-level
collaborative priors and are prone to hallucinations (e.g., sycophancy).
Naively integrating them often leads to negative transfer, where
semantic noise dilutes structural precision. To resolve this dilemma, we
propose ASC-Rec (Adaptive Structural-Semantic Cascade), a novel
framework operating on an "Expansion-Compression" paradigm.

In the Expansion phase, we construct a heterogeneous expert pool
spanning structural, statistical, and semantic views. An Attention-based
Dynamic Router acts as an adaptive gatekeeper, capturing diverse user
states (e.g., cold-start vs. active) to retrieve a broad, high-quality
candidate reservoir. In the Compression phase, we treat the LLM not as a
direct ranker, but as a Semantic Denoiser. We introduce a Conditional
Reciprocal Rank Fusion (Conditional RRF) strategy that leverages
LLM-generated reasoning to validate candidates. Crucially, this
mechanism applies "soft demotion" to items flagged as irrelevant,
suppressing noise without discarding valid structural signals. Extensive
experiments on three real-world datasets (Amazon Beauty, Sports, and
Toys) demonstrate that ASC-Rec significantly outperforms
state-of-the-art baselines. Notably, analysis of information density
reveals that our approach successfully shifts the "Golden Ratio" of
relevant items from the tail to the head of the ranking list, maximizing
recommendation utility within limited exposure slots.

\textbf{1. INTRODUCTION}

\textbf{Consumer decision-making is rarely driven by a singular factor.}
It is, in essence, a complex cognitive symphony orchestrated by a
multitude of entangled dimensions: the ripple effect of social trends
("What are my peers buying?"), the continuity of temporal habits ("What
fits my current lifestyle?"), and the intrinsic attraction to item
semantics ("Does this product's feature match my specific need?"). Even
for a cold-start user with limited interactions, every click is a
holographic projection of these underlying dimensions. Therefore, a
robust recommender system must act as a \textbf{polymath}, capable of
analyzing user intent from every possible angle to piece together a
complete user portrait.

However, prevalent recommendation paradigms often suffer from
"\textbf{Tunnel Vision}". Traditional Collaborative Filtering (CF)
methods, predominantly fixate on a single dimension---\textbf{structural
co-occurrence}---reducing rich human behaviors into abstract ID
connections. While efficient, this approach inevitably fails when the
structural signal is sparse or when the user's motivation is
content-driven (e.g., searching for a specific function) rather than
socially driven. Relying on such a "Flat Projection" of a
multi-dimensional world naturally leads to blind spots, resulting in
suboptimal recommendations, especially for long-tail items and
cold-start users.

To dismantle this barrier, we advocate for a paradigm shift from
"Single-View Modeling" to "\textbf{Multi-Perspective Perception}". We
propose \textbf{ASC-Rec} (Adaptive Structural-Semantic Cascade), a
framework designed to mimic the multi-dimensional nature of human
decision-making. Our approach is built upon two core philosophies:

\textbf{1. Expansion via Omni-Directional Experts:}

Just as a human decision is weighed against multiple factors, we
construct a Multi-View Expert Pool comprising six specialized agents,
each probing a distinct dimension of user intent:

\begin{itemize}
\item
  The Structural \& Sequence Experts (StructCLR, SASRec): Capture
  implicit habits and temporal dynamics.
\item
  The Statistical Experts (EASE, Tribe): Uncover global correlations and
  group-level trends (the "Wisdom of the Crowd").
\item
  The Semantic Experts (BM25, UserKNN): Decode explicit content matching
  and semantic taste similarities.
\end{itemize}

By employing an \textbf{Attention-based Dynamic Router}, our system
adaptively shifts its focus among these dimensions---much like how a
human might prioritize "brand" over "popularity" in different
contexts---thereby maximizing the breadth of the candidate reservoir.

\textbf{2. Compression via Cognitive Reasoning:}

While expanding perspectives increases recall, it also introduces
noise---recommendations that are statistically correlated but logically
flawed (e.g., recommending a phone case for a phone the user didn't
buy). To resolve this, we introduce a \textbf{Semantic Compression}
stage. We leverage the reasoning power of Large Language Models (LLMs)
not merely as a ranker, but as a "Cognitive Filter". Through a novel
Conditional Reciprocal Rank Fusion (Conditional RRF) mechanism, the LLM
scrutinizes the candidates, identifying and suppressing logically
irrelevant items ("No clear association"). This effectively
\textbf{compresses} the information entropy, distilling the chaotic
Top-50 candidates into a high-density, high-precision Top-K list.

In summary, this work makes the following contributions:

\begin{itemize}
\item
  Perspective:We reframe recommendation as a multi-dimensional
  perception task, proposing a MoE architecture that integrates
  Structure, Statistics, and Semantics.
\item
  Methodology:We develop the ASC-Rec framework, featuring an Adaptive
  Router for expansion and an LLM-driven Conditional RRF for semantic
  compression.
\item
  Empirical Results:Extensive experiments on Amazon Beauty, Sports, and
  Toys datasets demonstrate that ASC-Rec achieves a new
  state-of-the-art, effectively shifting the "Golden Ratio" of relevant
  items from the tail to the head of the ranking list.
\end{itemize}

\textbf{2. RELATED WORK}

In this section, we review the literature relevant to our work,
categorizing it into three streams: Graph-based Collaborative Filtering,
Large Language Models for Recommendation, and Mixture-of-Experts
mechanisms.

\textbf{2.1 Graph-based Collaborative Filtering}

Graph Neural Networks (GNNs) have advanced the modeling of high-order
user-item connectivity. Early methods like NCF and BPR-MF {[}Rendle et
al., 2009{]} focused on matrix factorization with pairwise ranking loss.
NGCF {[}Wang et al., 2019{]} later explicitly modeled collaborative
signals via message passing. LightGCN {[}He et al., 2020{]} simplified
the architecture by removing non-linear activations, establishing a
robust baseline. To address data sparsity, recent works like SimGCL
{[}Yu et al., 2022{]} integrated contrastive learning.

\textbf{Limitations}: Despite their success, these ID-based models
suffer from "Semantic Blindness." They rely solely on interaction
topology. Even advanced sequence models like BERT4Rec {[}Sun et al.,
2019{]}, which applies bidirectional Transformers to interaction
sequences, still operate primarily within the ID space. In contrast, our
framework employs StructCLR, an enhanced graph encoder with explicit
hard negative mining, as just one expert in a broader multi-view
ecosystem, ensuring structural signals are complemented by semantic
reasoning.

\textbf{2.2 Large Language Models in Recommendation}

The emergence of LLMs has sparked a paradigm shift from ID-based to
Content-based reasoning. Existing approaches generally fall into two
categories:

\begin{itemize}
\item
  \textbf{Tuning-based Methods:} These methods fine-tune open-source
  LLMs on interaction data to align them with recommendation tasks.
  TALLRec {[}Bao et al., 2023{]} proposes an efficient tuning framework
  to align LLMs with recommendation instructions using LoRA. CoLLM
  {[}Zhang et al., 2023{]} attempts to bridge the gap by distilling
  knowledge from a collaborative embedding model into a frozen LLM via a
  mapping network. PALR {[}Yang et al., 2023{]} further introduces
  cross-attention mechanisms to enhance user-item interaction modeling.
  While effective, these methods incur high training costs and face the
  risk of "Catastrophic Forgetting" of general world knowledge.
\item
  \textbf{Inference-based \& Agent Methods:} These methods utilize LLMs
  as zero-shot rankers or autonomous agents. GPT4Rec {[}Xu et al.,
  2024{]} leverages graph prompting to enable generative recommendation
  with pre-trained knowledge. More recently, agent-based models like
  RecMind {[}Wang et al., 2024{]} simulate user-agent interactions to
  explore hypothetical preferences in open-ended conversations.
\end{itemize}

\textbf{Limitations}: A critical, often overlooked issue in these
generative methods is the "Sycophancy" or hallucination problem, where
LLMs generate plausible justifications for irrelevant items due to a
lack of population-level collaborative priors. ASC-Rec differs by
treating the LLM not as a direct ranker or a conversational agent, but
as a Semantic Compressor. We introduce a Conditional RRF strategy to
explicitly mitigate LLM hallucinations via soft demotion, leveraging the
strengths of GPT4Rec-style reasoning while anchoring it with structural
stability.

\textbf{2.3 Mixture of Experts (MoE) and Adaptive Retrieval}

To handle diverse user intents, Multi-Interest Learning and MoE
architectures have gained attention.

\begin{itemize}
\item
  Multi-Interest Learning:~Models like~MIND~{[}Li et al., 2019{]}
  and~ComiRec~{[}Cen et al., 2020{]} generate multiple vectors per user
  to capture polysemous interests.
\item
  MoE Architectures:~In multi-task learning,~MMoE~{[}Ma et al., 2018{]}
  and~PLE~{[}Tang et al., 2020{]} use gating networks to distribute
  input signals to specific experts.
\end{itemize}

Limitations: Most existing MoE recommenders operate within a homogeneous
data view. Our Attention Router bridges this gap by learning "Expert
Prototypes" in the latent space, enabling adaptive switching between
structural collaborative filtering, statistical co-occurrence (e.g.,
EASE), and semantics earch based on the user's real-time state.

\hypertarget{methodology}{%
\subsubsection{3. METHODOLOGY}\label{methodology}}

\hypertarget{preliminaries}{%
\paragraph{3.1 Preliminaries}\label{preliminaries}}

In this section, we formally define the notation and the recommendation
problem addressed in this work.

Let \(\mathcal{U} = \{ u_{1},u_{2},\ldots,u_{M}\}\) be the set of M
users and \(\mathcal{V} = \{ v_{1},v_{2},\ldots,v_{N}\}\) be the set of
N items. We focus on the implicit feedback scenario, where user
interactions are represented by a binary matrix
\(Y \in \{ 0,1\}^{M \times N}\). An entry \(y_{\text{uv}} = 1\)
indicates that user \(u\) has interacted with item \(v\) (e.g.,
purchased or clicked), while \(y_{\text{uv}} = 0\) denotes unobserved
interactions. For each user \(u\) , we denote the set of interacted
items as \(\mathcal{H}_{u} = \{ v \in \mathcal{V}|y_{\text{uv}} = 1\}\)
.

To support our Multi-Expert architecture, we model the data from~three
distinct views, corresponding to the inputs required by our Structural,
Statistical, and Semantic experts:

\begin{itemize}
\item
  User-Item Graph: We construct a bipartite graph
  \(\mathcal{G} = \left( \mathcal{W},\mathcal{E} \right)\), where the
  node set \(\mathcal{W} = \mathcal{U} \cup \mathcal{V}\) includes all
  users and items, and the edge set \(\mathcal{E}\) corresponds to the
  observed interactions in \(Y\). This graph serves as the foundation
  for our structural experts (e.g., CL-GCN).
\item
  Textual Attributes: Each item \(v\) is associated with a textual
  sequence \(T_{v}\)(e.g., title, brand, description). This semantic
  modality enables our semantic experts and the LLM module to perform
  content-aware reasoning.
\item
  Graph View (Structural Signal):We construct a user-item bipartite
  graph \(\mathcal{G} = \left( \mathcal{W},\mathcal{E} \right)\), where
  the node set \(\mathcal{W} = \mathcal{U} \cup \mathcal{V}\) includes
  all users and items, and the edge set \(\mathcal{E}\) corresponds to
  the observed interactions in \(Y\). This graph serves as the
  foundation for structural experts (e.g., StructCLR) to perform message
  passing.
\item
  Co-occurrence View (Statistical Signal):Beyond local connectivity, we
  consider the global item- item correlation. This is implicitly modeled
  by the item similarity matrix derived from
  \(\mathbf{Y}^{T}\mathbf{Y}\), which captures the statistical
  co-occurrence patterns between items across the entire population.
  This view empowers our statistical experts (e.g., EASE, Tribe).
\item
  Semantic View (Content Signal):Each item \(v\) is associated with a
  textual sequence \(T_{v}\) (e.g., title, brand, description). This
  modality enables our semantic experts (e.g., BM25, UserKNN) and the
  LLM module to perform content-aware reasoning and denoising.
\end{itemize}

\textbf{Problem Definition (Top-K Recommendation)}. Given the
interaction history \(Y\) and item textual attributes \(\{ T_{v}\}\) ,
our goal is to learn a scoring function \(f\left( u,v \right)\) that
predicts the preference score of user \(u\) for an unobserved item
\(v\). Based on these scores, we aim to distill a high-precision ranked
list of K items from the vast item space \(\mathcal{V}\). This requires
not only retrieving a broad set of relevant candidates (Expansion) but
also rigorously filtering out structural noise to maximize information
density in the top positions (Compression).

\begin{longtable}[]{@{}ll@{}}
\toprule
\endhead
Symbol & Description\tabularnewline
\(U,\mathcal{V}\) & Sets of users and items\tabularnewline
\(M,N\) & Number of users and items\tabularnewline
\(Y\) & User-item interaction matrix\tabularnewline
\(\mathcal{G}\) & User-item bipartite graph\tabularnewline
\(E_{u},E_{v}\) & Learned embeddings for user \(u\) and item
\(v\)\tabularnewline
\(T_{v}\) & Textual content of item \(v\)\tabularnewline
\(\mathcal{C}_{u}\) & The candidate item set retrieved by
experts\tabularnewline
\(W_{\text{pos}},W_{\text{neg}}\) & Fusion weights for conditional
RRF\tabularnewline
\bottomrule
\end{longtable}

\hypertarget{framework-overview}{%
\paragraph{3.2 Framework Overview}\label{framework-overview}}

To transcend the inherent limitations of single-view representations and
effectively bridge the gap between collaborative signals and semantic
reasoning, we propose ASC-Rec (Adaptive Structural-Semantic Cascade). As
illustrated in Figure 1, the proposed framework operates under an
"Expansion-Compression" paradigm, treating the recommendation process as
a progressive information distillation pipeline consisting of three
stages:

\includegraphics[width=5.76528in,height=2.88264in]{./media/image1.png}

1. (Candidate Reservoir Construction):

Recognizing that user intent is multifaceted and dynamic, we posit that
relying on a single modality leads to suboptimal coverage. To address
this, we construct a heterogeneous Expert Pool encompassing three
distinct views: Structure, Statistics, and Semantics. This stage aims to
maximize the Recall Upper Bound by aggregating diverse signals---ranging
from high-order connectivity to keyword matching---thereby ensuring that
the initial candidate pool covers the full spectrum of potential user
interests.

2. Adaptive Routing (Structural Selection):

Simple aggregation of experts often introduces noise and redundancy. To
mitigate this, we introduce an Attention-based Dynamic Router. Acting as
an adaptive gating mechanism, the Router computes user-specific
importance weights based on the user's structural embedding. This stage
functions as a "Filter", efficiently retrieving a high-quality candidate
subset (e.g., Top-50) that aligns with the user's dominant behavioral
pattern, effectively establishing a robust structural prior.

3. Semantic Compression (Cognitive Refinement):

While the Router ensures high recall, the candidate list inevitably
contains structurally plausible but semantically irrelevant noise (e.g.,
false positives due to popularity bias). In the final stage, we employ a
Large Language Model (LLM) as a Semantic Compressor. By leveraging a
novel Conditional Reciprocal Rank Fusion (Conditional RRF) strategy, the
model incorporates semantic reasoning to identify and suppress negative
candidates ("No clear association"). This process compresses the entropy
of the ranking list, distilling valid items from the broad candidate
pool into the ultra-head ranking positions (Top-K), thereby maximizing
the utility within limited exposure slots.

\hypertarget{multi-view-expert-pool-expansion-stage}{%
\paragraph{3.3 Multi-View Expert Pool (Expansion
Stage)}\label{multi-view-expert-pool-expansion-stage}}

To ensure the initial candidate reservoir captures the full spectrum of
user intent, we construct a heterogeneous expert pool \(\mathcal{E}\).
The design philosophy is governed by Mutual Complementarity: experts
operate on orthogonal data views, ensuring that the scoring function
\(S_{k}\left( u,v \right)\) of the \(k\)-th expert covers the blind
spots of others.

3.3.1 Capturing Implicit Habits: The Structural View

Graph-based and sequence-based models serve as the backbone of our
system, excelling at capturing implicit user preferences from historical
interactions \(H_{u}\).

\begin{itemize}
\item
  \textbf{StructCLR} (Structural Expert)\textbf{:}
\end{itemize}

Standard GCNs often suffer from noise and uniform embedding
distribution. We propose StructCLR, an enhanced graph encoder. Unlike
vanilla LightGCN, StructCLR introduces a temperature-scaled InfoNCE loss
with explicit hard negative mining. By treating low-rated items (Rating
\textless{} 4) as hard negatives \(\mathcal{N}_{\text{hard}}\), we force
the model to distinguish fine-grained preferences. The objective is
defined as:

\[\mathcal{L}_{\text{Struct}} = - \sum_{u \in \mathcal{U}}^{}\log\frac{\exp\left( e_{u} \cdot e_{v^{+}}/\tau \right)}{\sum_{i \in \mathcal{V}}^{}\exp\left( e_{u} \cdot e_{i}/\tau \right) + \beta\sum_{j \in \mathcal{N}_{\text{hard}}}^{}\exp\left( e_{u} \cdot e_{j}/\tau \right)}\]

where \(\beta\) controls the penalty for confirmed negatives. The
prediction score is computed as the inner
product:\(S_{\text{struct}}\left( u,v \right) = e_{u}^{\top}e_{v}\).This
provides a robust structural prior for the router.

\begin{itemize}
\item
  \textbf{SASRec} (Sequence Expert)\textbf{:}
\end{itemize}

To capture temporal dynamics, SASRec models the interaction sequence
using self-attention. Given the history sequence~\(H_{u}\), the
prediction is based on the matching between the sequence
representation~\(h_{L}\)~and item embedding:

\[S_{\text{seq}}\left( u,v \right) = h_{L}\left( H_{u} \right)^{\top}e_{v}\]

3.3.2 Bridging Structural Gaps: The Statistical View

Structural models struggle to connect items that are topologically
distant or disjoint in the graph. We introduce statistical experts to
bridge these structural gaps via global co-occurrence patterns.

\begin{itemize}
\item
  \textbf{EASE} (Co-occurrence Expert)\textbf{:}
\end{itemize}

While GCN propagates information locally, EASE captures the global
item-item correlation matrix directly. It excels at identifying "hard
correlations" (e.g., functional complements) that embedding-based models
might miss due to the over-smoothing problem.

EASE computes a closed-form weight
matrix~\(W \in \mathbb{R}^{N \times N}\)~directly from the interaction
matrix~\(Y\). It captures global item-item correlations by solving a
constrained least-squares problem:

\[\min_{\mathbf{W}}\left| \  \right|\mathbf{Y} - \mathbf{\text{YW}}\left| \  \right|_{F}^{2} + \lambda\left| \  \right|\mathbf{W}\left| \  \right|_{F}^{2},\quad\text{s.t.\:diag}\left( \mathbf{W} \right) = 0\]

The score is computed via sparse matrix
multiplication:~\(S_{\text{ease}}\left( u,: \right) = y_{u} \cdot W\)

\begin{itemize}
\item
  \textbf{Tribe} (Cluster Expert)\textbf{:}
\end{itemize}

To mitigate the noise from individual behaviors, we aggregate users into
semantic clusters. This expert retrieves the most popular items within
each "Tribe," serving as a robust safety net against sparsity and
overfitting.This expert clusters users
into~\(K\)~groups~\(\{ C_{1},\ldots,C_{K}\}\). For a
user~\(u \in C_{k}\), the score is proportional to the item's popularity
within the cluster:

\[S_{\text{tribe}}\left( u,v \right) = \frac{\text{Count}\left( v \in C_{k} \right)}{\sum_{v^{\prime}}^{}\text{Count}\left( v^{\prime} \in C_{k} \right)}\]

3.3.3 Handling Explicit \& Cold-Start Intent: The Semantic View

Both structural and statistical models fail when interaction data is
scarce (Cold Start) or when user intent is explicitly content-driven. We
address this by incorporating semantic experts.

\begin{itemize}
\item
  \textbf{BM25 ~(Search Expert):}
\end{itemize}

Embedding models often suffer from "semantic drift." BM25 ensures
precision by performing rigid keyword matching between user history and
item titles, capturing explicit search-like intents (e.g., specific
brand or model requirements).We treat the user's history~\(H_{u}\)~as a
query and item title~\(T_{v}\)~as a document. The score represents the
keyword matching degree:

\[S_{bm25}\left( u,v \right) = \sum_{t \in \mathcal{H}_{u} \cap T_{v}}^{}\text{IDF}\left( t \right) \cdot \frac{\text{TF}\left( t,T_{v} \right) \cdot \left( k_{1} + 1 \right)}{\text{TF}\left( t,T_{v} \right) + k_{1} \cdot \left( 1 - b + b \cdot \frac{\left| T_{v} \right|}{\text{avgdl}} \right)}\]

\begin{itemize}
\item
  \textbf{UserKNN (Semantic Expert):}
\end{itemize}

This expert overcomes the "Product Disjoint" problem. By locating users
with similar semantic profiles (averaged BERT embeddings) rather than
overlapping purchase histories, it retrieves candidates from "semantic
soulmates," effectively handling scenarios where collaborative signals
are weak.To handle disjoint item sets, we represent each user as the
centroid of their historical item BERT
embeddings:~\(c_{u} = \frac{1}{\left| \mathcal{H}_{u} \right|}\sum_{i \in \mathcal{H}_{u}}^{}\text{BERT}\left( T_{i} \right)\).
The score is aggregated from the top-K~semantic
neighbors~\(\mathcal{N}_{u}^{\text{sem}}\):

\[S_{\text{knn}}\left( u,v \right) = \sum_{u^{\prime} \in \mathcal{N}_{u}^{\text{sem}}}^{}\text{Sim}\left( \mathbf{c}_{u},\mathbf{c}_{u^{\prime}} \right) \cdot \mathbb{I}\left( v \in \mathcal{H}_{u^{\prime}} \right)\]

\hypertarget{attention-based-dynamic-router}{%
\subsubsection{3.4 Attention-based Dynamic
Router}\label{attention-based-dynamic-router}}

Given the heterogeneous expert pool
\(\mathcal{E} = \{ E_{1},\ldots,E_{M}\}\) , a static aggregation
strategy (e.g., averaging) is suboptimal because different users exhibit
distinct behavioral patterns. For instance, a cold-start user may rely
heavily on the BM25 (Search) expert, while a community-driven user
aligns better with StructCLR (Structure).

To address this, we propose an Attention-based Dynamic Router that acts
as an adaptive traffic controller. It learns to map a user's current
state to the optimal expert combination.

\hypertarget{expert-prototypes-and-attention-mechanism}{%
\paragraph{3.4.1 Expert Prototypes and Attention
Mechanism}\label{expert-prototypes-and-attention-mechanism}}

We posit that each expert \(E_{k}\) specializes in a specific subspace
of user preferences. We learn a set of Expert Prototypes
\(P \in \mathbb{R}^{M \times d}\), where \(p_{k}\) represents the latent
centroid of users who are best served by expert \(E_{k}\).We utilize the
user embedding \(e_{u}\) from the pre-trained StructCLR as the query
vector, as it encodes the fundamental structural state of the user. The
attention weight \(\alpha_{u,k}\), representing the confidence assigned
to expert k for user u, is computed via a dot-product attention
mechanism:

\[\alpha_{u,k} = \frac{\exp\left( \mathbf{e}_{u}^{\top}\mathbf{p}_{k}/\tau_{\text{gate}} \right)}{\sum_{j = 1}^{M}\exp\left( \mathbf{e}_{u}^{\top}\mathbf{p}_{j}/\tau_{\text{gate}} \right)}\]

where \(\mathcal{T}_{\text{gate}}\) is a temperature parameter to
control the sparsity of the selection. A highly trusted expert receives
a weight \(\alpha_{u,k} \rightarrow 1\) , while irrelevant experts are
suppressed.

\hypertarget{weighted-reservoir-construction}{%
\paragraph{3.4.2 Weighted Reservoir
Construction}\label{weighted-reservoir-construction}}

Once the weights
\(\mathcal{A}_{u} = \{\alpha_{u,1},\ldots,\alpha_{u,M}\}\) are obtained,
we generate the initial candidate reservoir \(\text{Cu}\) . To combine
the ranked lists from different experts (which may have vastly different
score scales), we employ a Weighted Borda Count strategy. The aggregated
score for an item \(v\) is calculated as:

\[S_{\text{router}}\left( u,v \right) = \sum_{k = 1}^{M}\alpha_{u,k} \cdot \frac{1}{\text{rank}_{k}\left( v \right) + 1}\]

Items are then sorted by \(S_{\text{router}}\), and the Top-N(e.g.,N=50)
are selected to form the Candidate Reservoir passed to the next stage.
This ensures that the reservoir is dominated by the most competent
experts for the specific user \(u\).

\hypertarget{optimization-with-oracle-supervision}{%
\paragraph{3.4.3 Optimization with Oracle
Supervision}\label{optimization-with-oracle-supervision}}

A key challenge in MoE is the "collapse problem" where the gate
converges to a single expert. To prevent this, we employ a Supervised
Gating strategy rather than end-to-end RL.

Oracle Label Generation: For each user \(u\) in the validation set, we
retrospectively evaluate which expert successfully retrieved the
ground-truth item within the top ranks. We assign a "best expert" label
\(y_{u}^{\ast}\):

\[y_{u}^{\ast} = arg\max_{k}\mathbb{I}\left( v_{\text{truth}} \in TopK\left( E_{k} \right) \right)\]

Training: The Router is trained to predict this oracle label by
minimizing the Cross-Entropy loss:

\[\mathcal{L}_{\text{Router}} = - \sum_{u}^{}{\sum_{k = 1}^{M}\mathbb{I}}\left( y_{u}^{\ast} = k \right) \cdot log\left( \alpha_{u,k} \right)\]

This explicitly forces the Router to learn the correlation between user
embeddings (state) and expert capabilities (action).

3.5 LLM-driven Semantic Compression

While the Router effectively maximizes the recall upper bound
(Expansion), the candidate reservoir \(C_{u}\) inevitably contains
"Retrieval Noise"---candidates that are retrieved due to statistical
correlations or keyword overlaps but are functionally irrelevant to the
user's specific intent (e.g., false positives caused by popularity
bias).

To address this, we propose a Semantic Compression module. Unlike
traditional reranking, this module functions as a "Denoising Filter". It
leverages the reasoning capability of LLMs to identify these false
positives and applies a Conditional Reciprocal Rank Fusion (Conditional
RRF) strategy to suppress noise without discarding valid signals
established by the Router.

3.5.1 Context-Aware Reasoning Generation

Directly asking LLMs to output a ranking score often leads to poor
calibration. Instead, we adopt a "Generate-then-Discriminate" paradigm.
We employ GLM-4.6 as a Reasoning Generator to produce a textual
rationale \(R_{\text{uv}}\) for each user-item pair.

To ground the LLM's reasoning and mitigate hallucinations, we inject
"Source Context" into the prompt. We explicitly inform the LLM which
expert retrieved the item (e.g., "Matched by BM25" or "Recommended by
StructCLR").

\begin{itemize}
\item
  Prompting Strategy: We utilize Few-Shot Chain-of-Thought (CoT)
  prompting. The LLM is instructed to act as a strict auditor.
\item
  Negative Detection: Crucially, we instruct the LLM to output a
  specific flag token (e.g., "No clear association") if the item is
  semantically incompatible with the user's history.
\end{itemize}

3.5.2 Discriminative Semantic Scoring

To quantify the generated rationales, we utilize a fine-tuned
Cross-Encoder (BGE-Reranker) as the Discriminator. The discriminator
computes a semantic relevance score \(S_{\text{sem}}\left( u,v \right)\)
by jointly encoding the user history \(H_{u}\) and the item profile
augmented with the generated rationale:

\[S_{\text{sem}}\left( u,v \right) = \sigma\left( \mathrm{\text{CrossEncoder}}\left( H_{u},T_{v} \oplus \left\lbrack \text{SEP} \right\rbrack \oplus R_{\text{uv}} \right) \right)\]

where \(\oplus\) denotes concatenation. By conditioning the score on the
generated rationale \(R_{\text{uv}}\), the discriminator can distinguish
between "strong recommendations" and "weak recommendations".

3.5.3 Conditional RRF with Soft Demotion

A critical challenge in LLM-based recommendation is the False Negative
problem. Our empirical analysis reveals that LLMs may erroneously reject
implicitly correlated items (e.g., functional complements found by EASE)
due to a lack of collaborative knowledge.

To solve this, we propose a Conditional RRF strategy with an Asymmetric
Weighting Mechanism. We fuse the base rank from the Router
\(\left( \text{ran}k_{\text{router}} \right)\) and the semantic rank
from the Cross-Encoder \(\left( \text{ran}k_{\text{sem}} \right)\) as
follows:

\[S_{\text{final}}\left( u,v \right) = \underset{\mathrm{\text{Router\ Prior}}}{\underbrace{\frac{1.0}{\eta + \text{ran}k_{\text{router}}}}} + \underset{\mathrm{\text{Semantit\ Adjustment}}}{\underbrace{\frac{\lambda}{\eta + \text{ran}k_{\text{sem}}}}}\]

where \(\eta\) is a smoothing constant (e.g., 60). The semantic weight
\(\lambda\) is dynamic, conditioned on the sentiment of the rationale
\(R_{\text{uv}}\):

\[\lambda = \left\{ \begin{matrix}
w_{\text{pos}} & & & \text{if\:}R_{\text{uv}}\mathrm{\text{is\ Positive}} \\
w_{\text{neg}} & \text{if\:}R_{\text{uv}}\mathrm{is\ Negative\ ("No\ clear\ association")} & & \\
\end{matrix} \right.\ \]

\begin{itemize}
\item
  Positive Boost (\(w_{\text{pos}}\)=0.5):When the LLM identifies a
  clear semantic match, we amplify its score to promote the item to the
  top.
\item
  Soft Demotion (\(w_{\text{neg}}\)=0.3):When the LLM flags an item as
  irrelevant, we assign a lower weight but do not zero it out. This
  implements a "Soft Demotion" strategy: if an item has a very high
  priority in the Router (indicating strong collaborative or statistical
  evidence), the robust Router Prior ensures it remains in the
  recommendation list, protecting against LLM misjudgment. Conversely,
  if an item is both statistically weak and semantically irrelevant, it
  naturally sinks to the bottom.
\end{itemize}

This mechanism effectively compresses the information entropy,
distilling the candidate list to a high-density state where top-ranked
items satisfy both multi-view retrieval criteria and semantic coherence.

\hypertarget{experiments}{%
\subsection{4. EXPERIMENTS}\label{experiments}}

\hypertarget{experimental-settings}{%
\subsubsection{4.1 Experimental Settings}\label{experimental-settings}}

\hypertarget{datasets}{%
\paragraph{4.1.1 Datasets}\label{datasets}}

We evaluate our proposed framework on three real-world benchmark
datasets derived from the Amazon Review Data (2023 version): Amazon
Beauty, Amazon Sports, and Amazon Toys. These datasets vary
significantly in domain, scale, and sparsity, providing a comprehensive
testbed for evaluating recommendation performance across different
scenarios.

\begin{itemize}
\item
  Beauty: A dataset with relatively dense interactions, serving as a
  standard testbed for full-population evaluation.
\end{itemize}

\begin{itemize}
\item
  Sports \& Toys: Large-scale datasets characterized by sparse
  interactions and long-tail distributions. These allow us to assess the
  model's capability in handling massive item spaces and extracting
  signals from sparse behaviors.
\end{itemize}

To ensure data quality, we convert explicit ratings ≥4.0 into positive
implicit feedback, while ratings \textless{}4.0 are retained as explicit
hard negatives for contrastive learning. We apply a recursive 3-core
filtering for both users and items. The statistics of the processed
datasets are summarized in Table 1.

Table 1: Statistics of the processed datasets

\begin{longtable}[]{@{}llll@{}}
\toprule
\endhead
Dataset & Users & Items & Interactions\tabularnewline
Beauty & 631986 & 115709 & 701528\tabularnewline
Sports & 6703391 & 957764 & 12980837\tabularnewline
Toys & 4204994 & 624792 & 8201231\tabularnewline
\bottomrule
\end{longtable}

\hypertarget{evaluation-protocols}{%
\paragraph{4.1.2 Evaluation Protocols}\label{evaluation-protocols}}

We adopt the All-Item Full Ranking protocol. For each test user, the
ground-truth item is ranked against all other non-interacted items in
the dataset.

\begin{itemize}
\item
  Metrics: We report Recall@K and NDCG@K with K∈\{5,10,20\}K
  \textbackslash{}in \textbackslash{}\{5, 10,
  20\textbackslash{}\}K∈\{5,10,20\}
\item
  Evaluation Strategy: For the Beauty dataset, we perform evaluation on
  the full test set. For the large-scale Sports and Toys datasets,
  considering the computational cost of LLM inference, we adopt a
  Representative Subsampling Strategy. We randomly sample a subset of
  users (N=1,000) that preserves the original data distribution
  (covering both head and tail users) to evaluate the semantic
  compression stage. Preliminary experiments confirm that this subset
  size is statistically sufficient to represent the model's performance
  on long-tail data.
\end{itemize}

\hypertarget{baselines}{%
\paragraph{4.1.3 Baselines}\label{baselines}}

We compare ASC-Rec with state-of-the-art methods from two categories:

1. Collaborative Filtering \& Graph Methods:

\begin{itemize}
\item
  NCF: A neural collaborative filtering framework that models non-linear
  user-item interactions.
\item
  ENMF: An efficient neural matrix factorization method that learns from
  all data without negative sampling.
\item
  LightGCN: The state-of-the-art graph model that simplifies GCN by
  focusing solely on neighbor aggregation.
\end{itemize}

2. LLM-enhanced Methods:

\begin{itemize}
\item
  CoRAL: A method that aligns LLM knowledge with collaborative signals
  via a shared embedding space.
\item
  ICL (In-Context Learning): A zero-shot approach where user behaviors
  are constructed as prompts for direct LLM prediction.
\item
  ReAct: An agent-based approach where the LLM iteratively reasons and
  searches to generate recommendations.
\end{itemize}

3. Proposed Method:

\begin{itemize}
\item
  ASC-Rec: Our proposed framework. Unlike conventional models, ASC-Rec
  operates on an "Expansion-Compression" paradigm. It first utilizes a
  Multi-Expert Router to expand the recall scope
  (Structure/Statistics/Semantics) and then employs an LLM-driven
  Conditional RRF module to compress semantic noise, achieving adaptive
  structural-semantic alignment.
\end{itemize}

\hypertarget{implementation-details}{%
\paragraph{4.1.4 Implementation Details}\label{implementation-details}}

Expert Pool Configuration:\\
To construct a diverse candidate pool, we configure the six experts with
standardized settings. For the neural-based experts (CL-GCN and SASRec),
we fix the embedding dimension to 64 and optimize them using Adam.
Specifically, CL-GCN employs 3 graph propagation layers, while SASRec is
set with a maximum sequence length of 50 and 2 attention blocks. For the
statistical and retrieval-based experts, we set the regularization
parameter λ=200 for EASE. The UserKNN and Tribe (Clustering) experts are
configured to retrieve from the top-50 nearest neighbors or cluster
centroids, respectively, while BM25 utilizes standard text tokenization
on item titles without additional pre-training.

Router \& Compression Settings:\\
The Attention Router projects user embeddings into the expert prototype
space. For the semantic compression stage, we deploy GLM-4 as a Semantic
Denoiser to identify irrelevant candidates (labeled as "No clear
association"). Instead of relying on LLM for positive ranking, we adopt
a Negative Suppression Strategy in the final fusion. Specifically, items
identified as irrelevant by the LLM are assigned a penalty weight (or
demoted to the tail), while potentially relevant items retain their
structural ranking scores. This effectively compresses the candidate
space by filtering out semantic noise, allowing valid items hidden in
the tail (Top 20-50) to surface into the Top-20 window. All experiments
are conducted on a Linux server equipped with NVIDIA GPUs and PyTorch
3.10.

\hypertarget{overall-performance}{%
\subsection{4.2 Overall Performance}\label{overall-performance}}

We compare the performance of ASC-Rec with competitive baselines on
Amazon Beauty, Sports, and Toys datasets. The detailed experimental
results are reported in Tables 2, 3, and 4.

\hypertarget{table-2-performance-on-amazon-beauty-longtail}{%
\paragraph{Table 2: Performance on Amazon Beauty
（LongTail）}\label{table-2-performance-on-amazon-beauty-longtail}}

\begin{longtable}[]{@{}ll@{}}
\toprule
\endhead
Method & Beauty\tabularnewline
& Recall@5\tabularnewline
Traditional ID-based Methods\tabularnewline
NCF & 0.0053\tabularnewline
ENMF & 0.0138\tabularnewline
LightGCN & 0.0149\tabularnewline
Generative / LLM-based Methods\tabularnewline
ICL &\tabularnewline
ReACT &\tabularnewline
CoRAL & 0.0474\tabularnewline
Our Method\tabularnewline
Ours & 0.102\tabularnewline
\bottomrule
\end{longtable}

\hypertarget{table-3-performance-on-amazon-sports}{%
\paragraph{Table 3: Performance on Amazon Sports
}\label{table-3-performance-on-amazon-sports}}

\begin{longtable}[]{@{}ll@{}}
\toprule
\endhead
Method & Sports\tabularnewline
& Recall@5\tabularnewline
Traditional ID-based Methods\tabularnewline
NCF & 0.0469\tabularnewline
ENMF & 0.0474\tabularnewline
LightGCN & 0.0434\tabularnewline
Generative / LLM-based Methods\tabularnewline
ICL &\tabularnewline
ReACT &\tabularnewline
CoRAL & 0.0833\tabularnewline
Our Method\tabularnewline
Ours & 0.0823\tabularnewline
\bottomrule
\end{longtable}

\hypertarget{table-4-performance-on-amazon-toys}{%
\paragraph{Table 4: Performance on Amazon Toys
}\label{table-4-performance-on-amazon-toys}}

\begin{longtable}[]{@{}ll@{}}
\toprule
\endhead
Method & Toys\tabularnewline
& Recall@5\tabularnewline
Traditional ID-based Methods\tabularnewline
NCF & 0.0252\tabularnewline
ENMF & 0.0182\tabularnewline
LightGCN & 0.0317\tabularnewline
Generative / LLM-based Methods\tabularnewline
ICL &\tabularnewline
ReACT &\tabularnewline
CoRAL & 0.0798\tabularnewline
Our Method\tabularnewline
Ours & 0.0939\tabularnewline
\bottomrule
\end{longtable}

\hypertarget{fig_overall_performance_fixed_legend}{%
\subsubsection{\texorpdfstring{\protect\includegraphics[width=5.75556in,height=3.07083in]{./media/image2.png}}{fig\_overall\_performance\_fixed\_legend}}\label{fig_overall_performance_fixed_legend}}

\hypertarget{performance-analysis}{%
\subsubsection{4.2.1 Performance Analysis}\label{performance-analysis}}

1. Dominance over Single-View Models:\\
Structure-based methods (e.g., LightGCN) encounter a performance
bottleneck on sparse datasets (\emph{Sports} Recall@20 6.12\%) due to
the lack of high-order connectivity for long-tail items. Pure semantic
methods (e.g., ICL), while capable of zero-shot retrieval, struggle with
ranking precision due to the absence of collaborative signals. ASC-Rec
effectively bridges this gap, consistently outperforming all baselines
across three datasets, validating the necessity of a hybrid
architecture.

2. The Leap from Adaptive Routing:\\
The Attention Router serves as a powerful candidate generator. On
\emph{Sports}, the Router alone achieves a Recall@20 of 22.85\%, more
than doubling the performance of LightGCN. This significant margin
confirms that our MoE strategy---dynamically selecting experts based on
user embeddings---successfully captures diverse user intents that no
single expert could cover. This establishes a robust "Candidate
Reservoir" for the subsequent stage. In the context of our framework,
the Router successfully fulfills the "Expansion" objective, maximizing
the theoretical upper bound of recall.

3. Refinement via Semantic Compression:\\
Based on the strong Router baseline, ASC-Rec further pushes the
performance boundary. As shown in the tables, ASC-Rec improves the
Recall@20 to 23.06\% and Recall@5 to 10.73\%.\\
Crucially, this improvement is achieved in a context where the baseline
(Router) is already highly optimized. By employing the LLM as a Semantic
Denoiser, ASC-Rec successfully identifies false positives within the
candidate list. The Conditional RRF strategy ensures that while
structural signals (Router) determine the general ranking, semantic
signals (LLM) effectively rescue relevant items from the tail (Rank
20-50) and suppress irrelevant noise, resulting in a more refined and
accurate recommendation list.

\hypertarget{information-density-and-peak-shift}{%
\subsubsection{4.2.2 Information Density and Peak
Shift}\label{information-density-and-peak-shift}}

To illustrate the efficiency of our semantic compression module, we
analyze the Information Density distribution, defined as the recall gain
rate at each ranking position K. Figure 3 visualizes the cumulative
recall and density curves on the Sports dataset.

The "Early Saturation" Phenomenon:\\
As observed in Figure 3, the Router baseline exhibits a linear
accumulation pattern, indicating that valid items are scattered loosely
across the Top-50 list. In contrast, ASC-Rec demonstrates a distinct
Logarithmic Saturation pattern. The information density curve shows a
sharp Leftward Peak within the ultra-head interval
(K∈{[}1,5{]}).\includegraphics[width=2.75556in,height=1.76597in]{./media/image3.png}\includegraphics[width=2.75069in,height=1.76875in]{./media/image4.png}

Interpretation:\\
This phenomenon signifies a strategic Utility Concentration. The
LLM-driven module effectively compresses the "entropy" of the candidate
list, forcing the most relevant items---which might otherwise be buried
in lower positions (e.g., Rank 30)---to surface into the immediate view
(Rank 1-5).\\
This implies that ASC-Rec delivers a higher "Signal-to-Noise Ratio" in
the top-ranked positions. In real-world applications with limited
display slots (e.g., mobile feeds), this capability to maximize utility
within a minimal window size is often more valuable than broad retrieval
metrics.

\hypertarget{ablation-study-complete-version}{%
\subsubsection{4.3 Ablation Study (Complete
Version)}\label{ablation-study-complete-version}}

To verify the effectiveness of each component in ASC-Rec, we conduct
extensive ablation studies on both Amazon Beauty and Amazon Sports.

\hypertarget{effectiveness-of-adaptive-routing}{%
\paragraph{4.3.1 Effectiveness of Adaptive
Routing}\label{effectiveness-of-adaptive-routing}}

We first investigate the necessity of the Attention Router. We compare
our learned router against a static "Average Fusion" baseline, which
assigns equal weights （w=1/6 ) to all experts. We also list the
performance of the strongest single expert (LightGCN) for reference.

{[}Table 5: Impact of Routing Mechanism (Recall@20){]}\\
\textbar{} Dataset \textbar{} Best Single (LightGCN) \textbar{} Average
Fusion \textbar{} Attention Router \textbar{} Improvement (vs Avg)
\textbar{}\\
\textbar{} :-\/-\/- \textbar{} :-\/-\/-: \textbar{} :-\/-\/-: \textbar{}
:-\/-\/-: \textbar{} :-\/-\/-: \textbar{}\\
\textbar{} Amazon Beauty \textbar{} 0.1084 \textbar{} 0.0978 (↓)
\textbar{} 0.2285 \textbar{} +133.6\% \textbar{}\\
Analysis:\\
Table 5 reveals a critical insight into multi-expert fusion:

\begin{itemize}
\item
  The "Average" Trap: On the \emph{Beauty} dataset, simple Average
  Fusion (0.0978) actually performs worse than the single LightGCN model
  (0.1084). This is because weak experts (e.g., SASRec at 0.0383, EASE
  at 0.0457) introduce significant noise, diluting the high-quality
  signals from the structural expert.
\item
  The Router Solution: In contrast, the Attention Router achieves a
  Recall@20 of 0.2285, doubling the performance of LightGCN. This
  confirms that the Router effectively acts as a "Smart Gatekeeper". It
  learns to suppress weak experts when they are unreliable and activate
  them only when they provide complementary value (e.g., handling
  cold-start users via BM25), thus avoiding the negative transfer
  observed in static fusion.
\end{itemize}

\hypertarget{impact-of-semantic-compression-strategies}{%
\paragraph{4.3.2 Impact of Semantic Compression
Strategies}\label{impact-of-semantic-compression-strategies}}

This section validates the Conditional RRF strategy. We compare the
performance of different LLM utilization strategies on top of the Router
baseline.

\begin{longtable}[]{@{}ll@{}}
\toprule
\endhead
Strategy & Amazon Beauty\tabularnewline
Router Only & 0.2285\tabularnewline
Hard Filter & 0.2192 (↓ )\tabularnewline
Standard RRF & 0.2253 (↓)\tabularnewline
Conditional RRF & 0.2306 (↑)\tabularnewline
\bottomrule
\end{longtable}

Analysis:\\
The results in Table 6 demonstrate the necessity of our conditional
mechanism:

\begin{itemize}
\item
  Fragility of Hard Filtering: Directly removing items rejected by the
  LLM consistently hurts performance (dropping to
  \textasciitilde{}0.218). Error analysis shows that while LLMs are good
  at general reasoning, they suffer from a 16\% False Negative Rate on
  domain-specific implicit correlations (e.g., specific brand matches
  that look semantically unrelated).
\end{itemize}

\begin{itemize}
\item
  Effectiveness of Soft Demotion: Our Conditional RRF (wneg=0.3)
  successfully mitigates this issue. By "demoting" rather than
  "deleting" potential negatives, we allow the strong structural priors
  from the Router to rescue falsely rejected items. Meanwhile, the
  boosted weight for positive reasoning (wpos=0.5) helps valid long-tail
  items surface to the top, achieving the best overall performance
  (0.2306).
\end{itemize}

\hypertarget{conclusion-and-future-work}{%
\paragraph{5. CONCLUSION AND FUTURE
WORK}\label{conclusion-and-future-work}}

In this work, we presented ASC-Rec, a novel cascaded framework designed
to bridge the chasm between structural collaborative filtering and
semantic reasoning. Recognizing that neither ID-based models nor LLMs
can singly capture the full spectrum of user intent, we introduced an
"Expansion-Compression" paradigm. By constructing a multi-view expert
pool and employing an Attention-based Router, we successfully expanded
the retrieval boundary, capturing diverse user needs from structural
habits to explicit searches.

Crucially, we identified a critical limitation in applying LLMs for
recommendation: the propensity for hallucinations and false negatives,
which we termed "Sycophancy" and "Collaborative Blindness." To address
this, we transformed the role of the LLM from a direct ranker to a
Semantic Compressor. Our proposed Conditional RRF strategy---featuring a
"Soft Demotion" mechanism---successfully suppresses semantic noise
without discarding valid structural priors. Extensive experiments on
Amazon Beauty, Sports, and Toys datasets demonstrate that ASC-Rec not
only achieves state-of-the-art performance but also effectively shifts
the "Golden Ratio" of relevant items from the tail to the head of the
ranking list, maximizing information density in limited exposure slots.

Future Work. We plan to extend this research in three directions:

1. End-to-End Fine-tuning: Currently, the LLM functions as a frozen
inference engine. We aim to explore parameter-efficient fine-tuning
(e.g., LoRA) to align the LLM's reasoning more closely with
domain-specific collaborative signals.

2. Multimodal Expansion: Incorporating visual signals (product images)
into the expert pool to address cases where text descriptions are
ambiguous.

3. Reinforcement Learning for Routing: Replacing the supervised router
with a Reinforcement Learning agent to optimize for long-term user
satisfaction and diversity beyond immediate click-through rates.

\end{document}
